\chapter*{Glossaire et liste des abréviations}
\addcontentsline{toc}{chapter}{Glossaire et liste des abréviations}
\markboth{Glossaire et abréviations}{Glossaire et abréviations}

\section*{Sigles et abréviations}

\begin{longtable}{@{}p{2.8cm} p{11.5cm}@{}}
\toprule
\textbf{Sigle} & \textbf{Signification} \\
\midrule
\endfirsthead
\toprule
\textbf{Sigle} & \textbf{Signification} \\
\midrule
\endhead
\bottomrule
\endfoot
BFD & \textit{Block Flow Diagram} --- schéma de procédé par blocs \\
BFS & \textit{Bankable Feasibility Study} --- étude de faisabilité bancable \\
CAPEX & \textit{Capital Expenditure} --- dépenses d'investissement \\
CIL / CIP & \textit{Carbon in Leach} / \textit{Carbon in Pulp} --- procédés d'adsorption de l'or sur charbon actif \\
EPCM & \textit{Engineering, Procurement and Construction Management} \\
HSE & \textit{Health, Safety, Environment} --- hygiène, sécurité, environnement \\
IFC & \textit{Industry Foundation Classes} --- format d'échange des maquettes numériques \\
KPI & \textit{Key Performance Indicator} --- indicateur clé de performance \\
OPEX & \textit{Operating Expenditure} --- dépenses d'exploitation \\
P\&ID & \textit{Piping and Instrumentation Diagram} --- schéma de tuyauterie et d'instrumentation \\
PFD & \textit{Process Flow Diagram} --- schéma de procédé \\
PMC & \textit{Project Management Consultancy} --- assistance à la gestion de projet \\
R\&D & Recherche et Développement \\
ROI & \textit{Return On Investment} --- retour sur investissement \\
ROM & \textit{Run-Of-Mine} --- minerai brut, tel qu'extrait de la mine \\
SID & \textit{Safety In Design} --- sécurité intégrée dès la conception \\
SOP & \textit{Start Of Plant} --- date de démarrage de l'usine \\
\end{longtable}

\section*{Glossaire}

\begin{description}[style=nextline,leftmargin=0pt]
\item[Flowsheet (schéma de procédé)] Représentation de l'enchaînement des opérations unitaires permettant de transformer un minerai brut en produit valorisable (concentré, métal).
\item[Étude d'orientation \textit{(scoping study)}] Première évaluation économique et technique d'un projet, servant à décider de sa poursuite (précision indicative).
\item[Étude de pré-faisabilité] Étude approfondissant plusieurs scénarios techniques et estimant les coûts avec une précision intermédiaire.
\item[Étude de faisabilité] Étude détaillée validant la solution retenue, avec une estimation des coûts de l'ordre de $\pm$\,15--30\,\%, base de la décision d'investissement.
\item[Ingénierie de base \textit{(basic engineering)}] Dimensionnement des principaux équipements et des grands choix techniques du procédé.
\item[Ingénierie de détail \textit{(detailed engineering)}] Production de l'ensemble des plans et documents nécessaires à la construction.
\item[Convoyeur ROM] Convoyeur à bande transportant le minerai brut (\textit{run-of-mine}) depuis la mine vers l'unité de traitement.
\item[Pré-assemblage] Assemblage au sol de sous-ensembles avant leur levage et leur montage en position définitive.
\end{description}
